\begin{vidu}%[3P3-15.4G][Loai2] 
	Điện năng được truyền tải từ nơi cấp điện đến nơi tiêu thụ thông qua đường dây một pha có điện trở $R$. Biết rằng: công suất phát không thay đổi là $\mathcal{P}$, điện áp và dòng điện cùng pha nhau và nhu cầu sử dụng điện năng tại nơi tiêu thụ là $P_t$. Nơi tiêu thụ có dùng một máy hạ áp với đầu ra có điện áp hiệu dụng không đổi bằng $U_2=220\ V$. Khi điện áp phát là U thì tỉ số hạ áp nơi tiêu thụ là $k_1=42$ và đáp ứng được nhu cầu sử dụng điện năng là $\dfrac{21}{25}$. Nếu muốn cung cấp đủ điện năng cho nơi tiêu thụ với điện áp phát là $2U$ thì tỉ số hạ áp nơi tiêu thụ $k_{2}$ bằng
	\choice
	{56}
	{81}
	{78}
	{\True 100}
	\loigiai{
		\begin{longtable}{|p{0.17\textwidth}|p{0.17\textwidth}|p{0.17\textwidth}|p{0.17\textwidth}|p{0.17\textwidth}|}
			\hline %========================
			$P=UI$ &
			$U$ &
			$I$ &
			$P_t=U'I$ & $U'=kU_0$
			\\
			\hline %========================
			$\dfrac{U}{440}\left({3}\right)$ &
			$U\left({1}\right)$ &
			$\dfrac{1}{440}\left({2}\right)$ &
			$21\left({1}\right)$ & $42.220(1)$
			\\
			\hline %========================
			$\dfrac{U}{440}\left({4}\right)$ &
			$2U\left({1}\right)$ &
			$\dfrac{1}{880}\left({5}\right)$ &
			$25\left({1}\right)$ &$k_2=220(6)$
			\\
			\hline %========================
	\end{longtable}}
\end{vidu}
\begin{vidu}%[3P3-15.4G][Loai2] 
	Từ một trạm điện, điện năng được truyền tải đến nơi tiêu thụ bằng đường dây tải điện một pha, biết công suất truyền đến nơi tiêu thụ luôn không đổi, điện áp và cường độ dòng điện luôn cùng pha. Ban đầu, nếu trạm điện chưa sử dụng máy tăng áp và ở cuối đường dây dùng máy hạ áp lí tưởng có tỉ số vòng dây bằng 5 thì độ giảm điện áp trên đường dây tải điện bằng 0,1 lần điện áp hiệu dụng hai đầu tải tiêu thụ. Để công suất hao phí trên đường dây tải điện giảm 100 lần so với mức ban đầu thì ở trạm điện cần sử dụng máy tăng áp lí tưởng có tỉ số vòng dây cuộn thứ cấp so với cuộn sơ cấp \textbf{gần nhất} với giá trị nào sau đây?
	\choice
	{10}
	{\True 9,8}
	{9}
	{8,1}
	\loigiai{
		\begin{longtable}{|p{0.1\textwidth}|p{0.1\textwidth}|p{0.1\textwidth}|p{0.1\textwidth}|p{0.1\textwidth}|p{0.1\textwidth}|p{0.1\textwidth}|p{0.1\textwidth}|}
			\hline %========================
			$U_{SC1}$ &
			$U_{TC1}$ &
			$\Delta P$ &
			$\Delta U$ &
			I &
			$U_{SC2}$ &
			$U_{TC2}$ &
			$P_t=hs$\\
			\hline %========================
			$5,1.U_0$ &
			$5,1.U_0$ &
			$\Delta P=I^2R=I.\Delta U$ &
			$0,1. U_0$ &
			I &
			$5U_0$ &
			$U_0$ &
			$5U_0 I$\\
			\hline %========================
			$5,1.U_0$
			(*) &
			$50,01 U_0$ &
			$\dfrac{\Delta P}{100}$ &
			$0,01U_0$ &
			$\dfrac{I}{10}$ &
			$50U_0$ &
			&
			$5U_0 I$\\
			\hline %========================
		\end{longtable}
		Từ (*) và (**): hệ số tăng áp $k = 9,8$.
	}
\end{vidu}
\begin{vidu}%[3P3-15.4G][Loai2] 
	Điện năng được truyền tải từ nơi cấp điện đến nơi tiêu thụ thông qua một đường dây cố định có điện trở R. Biết điện áp nơi phát là $U_A=500\ V$ không đổi và nơi tiêu thụ đặt một máy ổn áp lí tưởng lấy điện đầu ra là $U_{2B}=200\ V$. Điện áp vào cuộn sơ cấp của ổn áp là không được nhỏ hơn 200 V và coi như điện áp cùng pha với cường độ dòng điện trên cả đường truyền. Khi nơi tiêu thụ cần sử dụng một công suất là 1,2 kW thì tỉ số ổn áp là 2. Khi nơi tiêu thụ cần sử dụng một công suất là 1100 W thì tỉ số ổn áp bây giờ \textbf{gần nhất} với giá trị nào sau đây?
	\choice
	{\True 2,05}
	{2,5}
	{1,45}
	{1,73}
	\loigiai{
		
		\begin{longtable}{|p{0.1\textwidth}|p{0.1\textwidth}|p{0.1\textwidth}|p{0.1\textwidth}|p{0.1\textwidth}|p{0.1\textwidth}|p{0.1\textwidth}|p{0.1\textwidth}|}
			\hline %========================
			$U$ &
			$\Delta U$ &
			$R$ &
			$I$ &
			$P_t$ &
			$U_t$ &
			$k$ &
			$U_t^{'}$\\
			\hline %========================
			$500(1)$ &
			$100(4)$ &
			$\dfrac{100}{3}(5)$ &
			$3(3)$ &
			$1200(1)$ &
			$400(2)$ &
			$2(1)$ &
			$200(1)$\\
			\hline %========================
			$500(1)$ &
			$500-200k(4)$ &
			$\dfrac{500-200k}{5,5k}(5)$ &
			$\dfrac{5,5}{k}(3)$ &
			$1100(1)$ &
			$200k(2)$ &
			$k(1)$ &
			$200(1)$\\
			\hline %========================
		\end{longtable}
		$(2)\Rightarrow U_t=kU;(3)\Rightarrow I=\dfrac{P_t}{U_t};(4)\Rightarrow \Delta U=U-U_t;(5)\Rightarrow R=\dfrac{\Delta U}{I}$
		$\dfrac{500-200k}{\dfrac{5,5}{k}}=\dfrac{100}{3}\Rightarrow k=2,05$.
	}
\end{vidu}
\begin{vidu}%[3P3-15.4G][Loai2] 
	\nguon{QG - 2020} immini[thm]{Điện năng được truyền tải từ máy hạ áp A đến máy hạ áp B bằng đường dây tải điện một pha như sơ đồ hình bên. Cuộn sơ cấp của A được nối với điện áp xoay chiều có giá trị hiệu dụng U không đổi, cuộn thứ cấp của B được nối với tải tiêu thụ X. Gọi tỉ số vòng dây của cuộn sơ cấp và số vòng dây của cuộn thứ cấp của A là $k_1$, tỉ số giữa số vòng dây của cuộn sơ cấp và số vòng dây của cuộn thứ cấp của B là $k_2$}
	{\begin{tikzpicture}[scale=1,thick,>=stealth']
		\draw[snake=coil,segment length=5pt,segment amplitude=6pt,line before snake=0.1cm,line after snake=0.1cm] (0,0) -- (0,3);
		\draw[snake=coil,segment length=5pt,segment amplitude=6pt,line before snake=0.5cm,line after snake=0.4cm] (0.75,0) -- (0.75,3);
		\draw(-0.5,1.25)--(-0.5,0)--(0,0)(0,3)--(-0.5,3)--(-0.5,1.75)(0.75,3)--(2,3)(0.75,0)--(2,0)(2.5,3)--(3.5,3)(2.5,0)--(3.5,0)(4.25,3)--(5,3)(4.25,0)--(5,0);
		\draw[dashed](2,3)--(2.5,3)(2,0)--(2.5,0);
		\draw[snake=coil,segment length=6pt,segment amplitude=6pt,line before snake=0.1cm,line after snake=0.1cm] (3.5,0) -- (3.5,3);
		\draw[snake=coil,segment length=4.5pt,segment amplitude=6pt,line before snake=0.6cm,line after snake=0.6cm] (4.25,0) -- (4.25,3);
		\draw[R=X] (5,3) -- (5,0);
		\draw[fill=white](-0.5,1.25)circle(1pt);
		\draw[fill=white](-0.5,1.75)circle(1pt);
		\draw[red,line width=2pt](0.375,0.5)--(0.375,2.5)node[above,black]{$A$};
		\draw[red,line width=2pt](3.875,0.5)--(3.875,2.5)node[above,black]{$B$};
		\draw[dashed](-0.4,-0.4)--(-0.4,3.4)--(1.2,3.4)--(1.2,-0.4)--cycle;
		\draw[dashed](3.1,-0.4)--(3.1,3.4)--(4.7,3.4)--(4.7,-0.4)--cycle;
		\node at (-0.6,1.5){$U$};
		\end{tikzpicture}}	
	Ở tải tiêu thụ, điện áp hiệu dụng như nhau, công suất tiêu thụ điện như nhau trong hai trường hợp $k_1=32$ và $k_2=68$ hoặc $k_1=14$ và $k_2=162$. Coi các máy hạ áp là lí tưởng, hệ số công suất của các mạch điện luôn bằng 1. Khi $k_1=32$ và $k_2=68$ thì tỉ số giữa công suất hao phí trên đường dây truyền tải và công suất ở tải tiêu thụ là
	\choice
	{0,009}
	{1,107}
	{0,019}
	{\True 0,052}
	\loigiai{
		\begin{itemize}
			\item 
			\item 
			\item 
			\item 
			\item 
			\item 
		\end{itemize}	
	}
\end{vidu}
